\documentclass[10pt,a4paper]{article}
\usepackage{fontspec}\setmainfont{Noto Serif CJK KR}\setmonofont{Noto Sans Mono CJK KR}
\usepackage[margin=22mm]{geometry}\usepackage{calc}\newcommand{\real}[1]{#1}\usepackage{longtable,booktabs,array,xurl,hyperref,fancyvrb,framed}
\providecommand{\tightlist}{\setlength{\itemsep}{0pt}\setlength{\parskip}{0pt}}
\newenvironment{Shaded}{}{}
\setlength{\parskip}{4pt}\setlength{\parindent}{0pt}
\begin{document}
\begin{center}{\LARGE\bfseries 서경 책 작업 인수인계 v12}\end{center}
\hypertarget{uxc11cuxacbd-uxcc45-uxc791uxc5c5-uxc778uxc218uxc778uxacc4-v12}{%
\section{서경 책 작업 인수인계 (v12)}\label{uxc11cuxacbd-uxcc45-uxc791uxc5c5-uxc778uxc218uxc778uxacc4-v12}}

책: 《고려 서경(西京)의 위치를 추정하다: 여덟 가지 방법으로 동아시아 고전 지명을 좌표로 옮기기》 (이광조, 과학과영성출판사). 이 문서만 읽고 바로 이어서 작업할 수 있도록 정리했다.

\hypertarget{uxd604uxc7ac-uxc0c1uxd0dc-v12}{%
\subsection{1. 현재 상태 (v12)}\label{uxd604uxc7ac-uxc0c1uxd0dc-v12}}

\begin{itemize}
\tightlist
\item
  빌드: 컬러 577쪽, 흑백 577쪽. 정의되지 않은 참조 0건, Overfull hbox 35건(대부분 코드 블록의 긴 줄).
\item
  완료: 5·6·7·8장 확장(각 +100,000자 안팎), 고문헌 실례 25건 삽입(가상 문장 아래에 추가만), 7장 표 22개 캡션 추가, 코드 강조 끄기.
\item
  산출물 위치: Downloads와 \texttt{Desktop\textbackslash{}daily3\textbackslash{}seogyeong-book\textbackslash{}pdf}의 \texttt{seogyeong-book-v12-color.pdf}, \texttt{seogyeong-book-v12-bw.pdf}.
\item
  소스: 컨테이너 \texttt{/home/claude/book}(세션이 끊기면 사라짐), 스테이징 \texttt{/mnt/user-data/outputs/seogyeong-book/}, 이박사님 컴퓨터 \texttt{Desktop\textbackslash{}daily3\textbackslash{}seogyeong-book}. DB: \texttt{placeloc-v12.sqlite}(모든 원고·도구·자료 포함).
\item
  git: v11 커밋 미확인. win11-01 PowerShell에서 \texttt{git\ add\ -A;\ git\ commit\ -m\ "seogyeong-book\ v12:\ 5-8장\ 확장,\ 고문헌\ 실례\ 삽입"}.
\end{itemize}

\hypertarget{uxc0c1uxc2dc-uxaddcuxce59-uxbc14uxafb8uxc9c0-uxb9d0-uxac83}{%
\subsection{2. 상시 규칙 (바꾸지 말 것)}\label{uxc0c1uxc2dc-uxaddcuxce59-uxbc14uxafb8uxc9c0-uxb9d0-uxac83}}

\begin{itemize}
\tightlist
\item
  근거만 쓴다: 지어낸 출처·숫자·인용 금지. 모르면 ``모름'', 원문을 못 읽었으면 ``원문 대조 필요''.
\item
  중학교 2학년이 따라 읽을 수 있게 친절하게 쓴다. 분량 제한 없음(분권 가능).
\item
  사례는 1차 자료 우선. 2·3차 자료는 ``비판적으로 읽기''에서만 쓰고 계산 입력에 쓰지 않는다. 가상 사례는 사료 문형을 본떠 만들고 ``가상''이라고 표시한다.
\item
  자료의 층위: 1층 원사료 인용, 2층 판독문·계측, 3층 연구자 해석. 3층은 증거로 쓰지 않는다.
\item
  B5 판형, 본문은 왼쪽 여백 20\%, 표와 그림은 전체 폭. 지도에는 강과 산을 표시한다. APA 7판. 한자어는 한국 한자음으로 쓰고 병음은 괄호 안에만 둔다.
\item
  이박사님이라고 부른다. 파일명에 언더바를 쓰지 않고 버전을 붙인다. 문서 산출물은 Downloads로 전달한다.
\item
  기존 원고 문장은 고치지 않고, 아래에 추가만 한다(삽입 후 삽입분을 빼면 원본과 같아야 한다).
\end{itemize}

\hypertarget{uxc6d0uxace0-uxad6cuxc131uxacfc-uxae00uxc790-uxc218-uxacf5uxbc31-uxd3ecuxd568}{%
\subsection{3. 원고 구성과 글자 수 (공백 포함)}\label{uxc6d0uxace0-uxad6cuxc131uxacfc-uxae00uxc790-uxc218-uxacf5uxbc31-uxd3ecuxd568}}

\begin{longtable}[]{@{}
  >{\raggedright\arraybackslash}p{(\columnwidth - 4\tabcolsep) * \real{0.3333}}
  >{\raggedright\arraybackslash}p{(\columnwidth - 4\tabcolsep) * \real{0.3333}}
  >{\raggedright\arraybackslash}p{(\columnwidth - 4\tabcolsep) * \real{0.3333}}@{}}
\toprule\noalign{}
\begin{minipage}[b]{\linewidth}\raggedright
파일
\end{minipage} & \begin{minipage}[b]{\linewidth}\raggedright
글자 수
\end{minipage} & \begin{minipage}[b]{\linewidth}\raggedright
비고
\end{minipage} \\
\midrule\noalign{}
\endhead
\bottomrule\noalign{}
\endlastfoot
00-intro \textasciitilde{} 04-db & 3,355 \textasciitilde{} 11,660 & 아직 확장 안 함 \\
05-csp 제약 전파 & 138,728 & 확장 완료 \\
06-georef 지오레퍼런싱 & 126,660 & 확장 완료 \\
07-bayes 베이즈 & 120,686 & 확장 완료 \\
08-lcp 최소비용경로 & 115,646 & 확장 완료 \\
09-graph \textasciitilde{} 12-cross & 12,500 \textasciitilde{} 17,129 & 아직 확장 안 함 \\
13-compare \textasciitilde{} 15-conclusion & 4,478 \textasciitilde{} 5,381 & 서경 실제 사료 없음, 확장 필요 \\
16-nwkorea & 12,586 & 증거 E1\textasciitilde E12, G1 적재됨 \\
17-liao, 18-shandong, 19-manchuria & 3,966 \textasciitilde{} 5,071 & 서경 사료 비어 있음 \\
A-math, B-glossary, C-dbspec, notes & 100 \textasciitilde{} 2,500 & 부록 \\
\end{longtable}

\texttt{원고/실례-삽입안.md}는 빌드 대상이 아니다(삽입 초안 보관용).

\hypertarget{uxbe4cuxb4dc-uxbc29uxbc95}{%
\subsection{4. 빌드 방법}\label{uxbe4cuxb4dc-uxbc29uxbc95}}

\begin{verbatim}
cd /home/claude/book
python3 tools/worked.py          # 완전 풀이 예제 블록 생성(멱등)
python3 tools/spectables.py      # 지도 제원표 삽입(멱등)
python3 tools/addcaps.py         # 캡션 없는 표에 캡션 추가(멱등)
BOOK_VERSION=v13 python3 tools/build.py color
BOOK_VERSION=v13 python3 tools/build.py bw
BOOK_VERSION=v13 ASSET_DIR=<산출물폴더> python3 tools/mk-db.py build/placeloc-v13.sqlite
\end{verbatim}

\begin{itemize}
\tightlist
\item
  원고 매크로: \texttt{\{\{key:fmt\}\}}, \texttt{@@fig\ 경로\textbar{}라벨\textbar{}캡션@@}, \texttt{@@tab\ csv\textbar{}라벨\textbar{}캡션\textbar{}colspec\textbar{}note@@}, \texttt{@@refs\ chNN@@}. 표는 \texttt{Table:\ 캡션} + 빈 줄 + 파이프 표.
\item
  pandoc 수식: 닫는 \texttt{\$} 바로 뒤에 숫자가 오면 안 된다.
\item
  \texttt{tools/insert-real.py}는 실례 삽입(마커 \texttt{\textless{}!-\/-\ real-example:N:장:길이\ -\/-\textgreater{}}, 멱등). 삽입분 제거 후 원본과 같은지 검사할 수 있다.
\item
  폰트: Noto Serif CJK KR, Noto Sans CJK KR.
\end{itemize}

\hypertarget{uxb3c4uxad6cuxc640-uxc790uxb8cc}{%
\subsection{5. 도구와 자료}\label{uxb3c4uxad6cuxc640-uxc790uxb8cc}}

\begin{itemize}
\tightlist
\item
  \texttt{tools/expand-guide.md}: 확장 집필 지침(구체적 사례 원칙, 1차 자료 우선). 새 장을 확장할 때 에이전트에게 먼저 읽히고, 대상 파일·목표 글자 수를 지정한다.
\item
  \texttt{data/evidence-seogyeong-v1.csv}: 증거 E1\textasciitilde E12, G1(layer 열 포함). \texttt{data/real-examples-v1.csv}: 고문헌 실례 30행. \texttt{data/mountains-v1.csv}, \texttt{data/rivers-v1.csv}: 지도 제원.
\item
  \texttt{tools/calc07/}, \texttt{tools/expand05/}: 7장·5장 계산 스크립트(재현용).
\item
  참고문헌: \texttt{참고문헌/chNN.txt}, 검증 현황 \texttt{참고문헌/verified-v1.csv}.
\item
  서버 작업은 PowerShell 한 줄 명령으로 한다: \texttt{ssh\ -J\ z440@100.90.3.79\ z840@192.168.0.84\ \textquotesingle{}...\textquotesingle{}} (작은따옴표 바깥, 안쪽에 따옴표 없이). 파일 회수는 \texttt{scp\ -J\ z440@100.90.3.79\ z840@192.168.0.84:/home/z840/X\ .\textbackslash{}data\textbackslash{}X} 후 기기 브리지로 읽는다. z840 \texttt{/data/safe/dball2/}에 nahf\_db(고구려통사 PDF), kyujanggak\_db, oldmap\_db, sagodb-v2.sqlite 등이 있다.
\end{itemize}

\hypertarget{uxd574uxc57c-uxd560-uxc77c-uxc6b0uxc120uxc21cuxc704-uxc21c}{%
\subsection{6. 해야 할 일 (우선순위 순)}\label{uxd574uxc57c-uxd560-uxc77c-uxc6b0uxc120uxc21cuxc704-uxc21c}}

\begin{enumerate}
\def\labelenumi{\arabic{enumi}.}
\tightlist
\item
  \textbf{확장 안 한 장}: 00\textasciitilde04, 09\textasciitilde15, 17\textasciitilde19, 부록을 같은 방식으로 확장한다(장당 +80,000자 이상 목표, 사례·Q\&A·연습문제·작업지 포함). 에이전트 4개를 동시에 돌리고 끝나면 보고받는다.
\item
  \textbf{1차 원문 확보}: 『고려사』 문종 35년(1081) 8월 신유조 서경 궁궐 기사(RE01), 『신증동국여지승람』 대성산 동북쪽(RE28), 『해동역사』 구룡산 위(RE29)의 한문 원문. z840의 kyujanggak\_db, sagodb-v2, oldmap\_db에서 찾는 PowerShell 한 줄 명령을 이박사님께 드려 실행하게 한다.
\item
  \textbf{E1, E3, E6, E9의 한자 원문 보강}: E3은 『삼국사기』 \texttt{秋七月\ 移居平壤東黄城\ 城在今西京東木覓山中}을 전사본에서 확인했다(판본 대조 필요). 증거표 상태 칸을 ``전사본 대조 일치''로 올린다. 각자성석 제4석을 증거표에 추가한다.
\item
  \textbf{13\textasciitilde15, 17\textasciitilde19장}: 실제 서경 사료가 확보되면 채운다. 지금은 비어 있다.
\item
  \textbf{지명 좌표}: 목멱산, 구룡산, 대성산(구룡산과 같은 산일 가능성), 안학궁, 주궁지 좌표. 1리의 길이와 척·리 비율은 출처가 없어 가정 구간(0.36\textasciitilde0.48 km)으로만 쓰고 있다.
\item
  \textbf{서경 결과 계산}: 사료와 좌표가 모이면 5\textasciitilde8장 방법으로 실제 서경 추정을 계산한다. 지금은 방법 설명용 가상 사례뿐이다.
\item
  \textbf{02-prior 문장 정리}: ``위치가 기록의 대상이 된 것은 조선 초''는 『삼국사기』 지리 4의 서경=평양 논증과 충돌하지 않게 표현을 좁힌다(삽입 4의 메모).
\item
  \textbf{분권}: 쪽수가 600쪽을 넘으면 방법편(00\textasciitilde08), 방법편II+비교(09\textasciitilde15), 지역 적용(16\textasciitilde19)+사료집으로 나누고 빌드가 권별 PDF를 만들도록 한다. 사료집 부록(E1\textasciitilde E12 원문)도 만든다.
\item
  \textbf{미검증 참고문헌}: Zadeh 1986 원문, ESA Copernicus DEM DOI, Adansonia 2025 DOI, 임기환 2007·기경량 2020·田中俊明 2004 서지, 한국어 서적류. \texttt{tools/check-cites.py}는 한국어 ``연도 미상'' 줄을 인식하지 못한다.
\item
  \textbf{Overfull hbox 35건}: 대부분 8장 코드 블록. 필요하면 코드 글꼴을 줄이거나 줄바꿈한다.
\end{enumerate}

\hypertarget{uxc8fcuxc758uxd560-uxc810}{%
\subsection{7. 주의할 점}\label{uxc8fcuxc758uxd560-uxc810}}

\begin{itemize}
\tightlist
\item
  에이전트는 사용자가 턴을 중단하면 취소된다. 에이전트를 띄운 뒤 곧바로 턴을 끝내야 끝까지 돈다.
\item
  컨테이너 셸은 패키지 저장소만 열린다. 위키문헌은 WebFetch로만 읽히고 요약 모델을 거치므로 같은 구절을 두 번 이상 읽어 일치한 것만 ``일치''로 적는다. 한국사데이터베이스는 robots로 막혀 있다.
\item
  한문 원문을 기억으로 인용하지 않는다. 연구서 번역만 있으면 ``원문 대조 필요''로 둔다.
\item
  서경의 실제 위치 결론은 아직 없다. 가상 사례의 숫자를 서경 결과로 읽히게 쓰지 않는다.
\end{itemize}

\hypertarget{uxc99duxac70-uxbaa9uxb85d-uxc694uxc57d-dataevidence-seogyeong-v1.csv}{%
\subsection{8. 증거 목록 요약 (data/evidence-seogyeong-v1.csv)}\label{uxc99duxac70-uxbaa9uxb85d-uxc694uxc57d-dataevidence-seogyeong-v1.csv}}

E1 고려사 문종 35년(서경에서 동서 10여 리 좌우 궁궐; 연구서 인용), E3 삼국사기 고국원왕 13년(평양 동쪽 황성, 목멱산), E5\textasciitilde E7 장안성 규모·면적, E8·E10\textasciitilde E12 장안성 각자성석(제2·4·5석 현존), E6 대성산 동북쪽, E9 구룡산 위, G1 현재 평양 좌표. 층위와 판독 상태는 CSV의 layer, status 열을 본다.

\end{document}
